\documentclass[11pt]{article}

\usepackage[a4paper,margin=1in]{geometry}
\usepackage[T1]{fontenc}
\usepackage[utf8]{inputenc}
\usepackage{lmodern}
\usepackage{microtype}
\usepackage{amsmath,amssymb,mathtools,bm}
\usepackage{graphicx}
\usepackage{booktabs}
\usepackage{enumitem}
\usepackage{xcolor}
\usepackage{hyperref}

\hypersetup{
  colorlinks=true,
  linkcolor=blue!60!black,
  citecolor=blue!60!black,
  urlcolor=blue!60!black
}

\newcommand{\SK}{S_K}
\newcommand{\chip}{\chi_{\rm t}}
\newcommand{\chistart}{\chi_{\rm start}}
\newcommand{\Acal}{\mathcal{A}}
\newcommand{\Ical}{\mathcal{I}}
\newcommand{\dd}{\,\mathrm{d}}
\newcommand{\e}{\mathrm{e}}
\newcommand{\lam}{\lambda}
\newcommand{\phipeak}{\phi_{\rm peak}}

\title{Onset location for hyperspherical Bessel functions\\[3pt]
\large Airy/action derivation and numerical calibration of practical start rules}
\author{}
\date{Updated May 2026}

\begin{document}
\maketitle

\section{Goal}

For numerical integration or sampling of regular hyperspherical Bessel functions
one often wants a start point before which the function is negligible compared
with its peak amplitude.  The target used here is not pointwise relative error,
but
\[
  \frac{|\phi_l(\chi)|}{\max_\chi |\phi_l(\chi)|} < \varepsilon,
\]
with practical values \(\varepsilon=10^{-5}\) and \(10^{-4}\).

In flat space a familiar rule for the onset of the spherical Bessel function is
\[
  x=\beta\chi \simeq l-5l^{1/3},
\]
for a roughly \(10^{-5}\)-of-peak start.  This note derives the curved-space
action generalisation and gives a numerical calibration, including open and
closed geometries up to \(l=4000\).

Throughout the calibrated formulae use the Langer order
\[
  \boxed{\lam=l+\frac12}
\]
unless otherwise stated.  Replacing \(\lam\) by
\(\ell=\sqrt{l(l+1)}\) gives a very similar large-\(l\) rule, but the Langer
form is the more natural one for Airy/Olver turning-point formulae.

\section{Equation, turning point, and notation}

The regular hyperspherical Bessel function satisfies
\[
\phi_l''(\chi)
+2\frac{\SK'(\chi)}{\SK(\chi)}\phi_l'(\chi)
+\left[\beta^2-K-\frac{l(l+1)}{\SK^2(\chi)}\right]\phi_l(\chi)=0,
\]
where
\[
\SK(\chi)=
\begin{cases}
\sin\chi, & K=1,\\
\chi, & K=0,\\
\sinh\chi, & K=-1.
\end{cases}
\]
Writing \(\tilde u=\SK\phi_l\) gives
\[
  \tilde u''(\chi)
  +\left[\beta^2-\frac{l(l+1)}{\SK^2(\chi)}\right]\tilde u(\chi)=0.
\]
The onset rule replaces \(l(l+1)\) by \(\lam^2\), giving the effective local
wavenumber
\[
  Q_K(\chi)=\beta^2-\frac{\lam^2}{\SK^2(\chi)}.
\]
The turning point is
\[
  Q_K(\chip)=0,
  \qquad
  \SK(\chip)=\frac{\lam}{\beta}.
\]
Thus
\[
\chip=
\begin{cases}
\arcsin(\lam/\beta), & K=1,\\
\lam/\beta, & K=0,\\
\operatorname{arsinh}(\lam/\beta), & K=-1.
\end{cases}
\]
For \(K=1\) the folded interval \(0\leq\chi\leq\pi/2\) is assumed.

Define also
\[
  r=\frac{\beta}{\lam}.
\]
For closed models \(K=1\) one has \(r>1\).  For open models \(K=-1\), all
positive \(r\) are possible, and the low-\(r\) regime is the only place where a
single flat-space constant \(c\simeq5\) noticeably mis-calibrates the onset.

\section{Flat-space action calibration}

In flat space \(\SK(\chi)=\chi\) and the exact solution is
\[
  \phi_l(\chi)=j_l(\beta\chi).
\]
Put
\[
  x=\beta\chi,
  \qquad
  y=\frac{x}{\lam}.
\]
The flat forbidden action from \(x\) to the turning point \(x=\lam\) is
\[
\Acal_0(y)
=\lam\int_y^1\left(\frac{1}{u^2}-1\right)^{1/2}\dd u.
\]
The elementary integral is
\[
  \Ical_0(y)
  =
  \log\left(\frac{1+\sqrt{1-y^2}}{y}\right)-\sqrt{1-y^2},
  \qquad y<1.
\]
The flat onset ansatz
\[
  x_{\rm start}=\lam-c\lam^{1/3}
\]
has
\[
  y_c=1-c\lam^{-2/3}
\]
and hence action level
\begin{equation}
\boxed{
T_{\lam,c}
=\lam\Ical_0(y_c)
=\lam\left[
\log\left(\frac{1+\sqrt{1-y_c^2}}{y_c}\right)-\sqrt{1-y_c^2}
\right].
}
\label{eq:T_exact}
\end{equation}
For large \(\lam\),
\[
  T_{\lam,c}=\frac{2\sqrt2}{3}c^{3/2}+O(\lam^{-2/3}).
\]
Thus \(c=5\) corresponds to the nearly \(l\)-independent action level
\[
  T\simeq \frac{2\sqrt2}{3}5^{3/2}\simeq 10.54.
\]
The action, rather than \(x\) itself, is the invariant quantity to preserve in
curved space.

\section{Curved-space action start rule}

For \(\chi<\chip\), define the forbidden-region action
\begin{equation}
\Acal_K(\chi)
=
\int_\chi^{\chip}
\left[
\frac{\lam^2}{\SK^2(s)}-\beta^2
\right]^{1/2}\dd s.
\label{eq:curved_action}
\end{equation}
The curved generalisation of the flat start rule is
\begin{equation}
\boxed{
\Acal_K(\chistart)=T_{\lam,c}.
}
\label{eq:action_rule}
\end{equation}
Equivalently,
\[
\Ical_K(\chistart)
=
\Ical_0(1-c\lam^{-2/3}),
\]
where
\[
\Ical_K(\chi)=
\int_\chi^{\chip}
\left[\frac{1}{\SK^2(s)}-r^2\right]^{1/2}\dd s,
\qquad r=\beta/\lam.
\]
This is the recommended general recipe.  It is a scalar root find in
\(0<\chi<\chip\), using either numerical quadrature or the closed analytic
actions already used in uniform Olver mappings.

\section{Explicit local formula}

The action rule has a simple explicit approximation from expanding about the
turning point.  Let
\[
  \Delta=\chip-\chi,
  \qquad
  \sigma_t=|\SK'(\chip)|=
  \sqrt{1-K\frac{\lam^2}{\beta^2}}.
\]
For \(\chi=\chip-\Delta\),
\[
  \SK(\chi)=\SK(\chip)-\sigma_t\Delta+O(\Delta^2),
\]
and therefore
\[
  \frac{\lam^2}{\SK^2(\chi)}-\beta^2
  =\frac{2\beta^3\sigma_t}{\lam}\Delta+O(\Delta^2).
\]
Hence
\[
\Acal_K(\chip-\Delta)
=\frac{2}{3}\left(\frac{2\beta^3\sigma_t}{\lam}\right)^{1/2}
\Delta^{3/2}+O(\Delta^{5/2}).
\]
Equating this to the asymptotic flat value
\((2\sqrt2/3)c^{3/2}\) gives
\begin{equation}
\boxed{
\chistart
\simeq
\chip-
\frac{c\lam^{1/3}}
{\beta\left(1-K\lam^2/\beta^2\right)^{1/6}}.
}
\label{eq:explicit_rule}
\end{equation}
In practice clamp to the lower boundary:
\begin{equation}
\boxed{
\chistart^{\rm fast}
=
\max\left[0,
\chip-
\frac{c\lam^{1/3}}
{\beta\left(1-K\lam^2/\beta^2\right)^{1/6}}
\right].
}
\label{eq:fast_rule}
\end{equation}
This is exact in form for the flat scaling and is useful as a cheap start guess.
The action-inverted rule, Eq.~\eqref{eq:action_rule}, is preferred near the
closed equator \(K=1\), \(\lam/\beta\simeq1\), or whenever the start location
matters sharply.

\section{Numerical calibration}

The calibration used exact numerical values of \(\phi_l\) obtained from a stable
recurrence: upward recurrence in the safe oscillatory region and Miller backward
recurrence in the tail, with a Gegenbauer logarithmic derivative for the closed
\(K=1\) tail.  The sweep covered open and closed cases up to \(l=4000\), with a
broad range of \(r=\beta/\lam\).

For each model the exact start point was defined by
\[
  |\phi_l(\chistart)|=\varepsilon\max_\chi |\phi_l(\chi)|,
\]
using the first crossing before the main peak.  This exact start was then
converted to an effective \(c_{\rm eff}\) through the action equation
\[
\Acal_K(\chi_{\varepsilon})=T_{\lam,c_{\rm eff}}.
\]
Large \(c_{\rm eff}\) means the rule must start earlier to reach the requested
threshold.

\subsection{Results for \texorpdfstring{\(10^{-5}\)}{1e-5}}

For \(\varepsilon=10^{-5}\), the flat/closed/moderate-open calibration is very
close to \(c=5\), but the open low-\(r\) tail requires a larger value.

\begin{center}
\begin{tabular}{lcc}
\toprule
Region & 95th percentile \(c_{\rm eff}\) & max observed \(c_{\rm eff}\) \\
\midrule
Closed \(K=+1\), all tested ratios & 4.924 & 4.933 \\
Open \(K=-1\), \(r\ge0.2\) & 4.933 & 4.945 \\
Open \(0.05\le r<0.2\) & 4.981 & 4.986 \\
Open \(0.02\le r<0.05\) & 5.113 & 5.118 \\
Open \(0.01\le r<0.02\) & 5.242 & 5.250 \\
Open \(0.005\le r<0.01\) & 5.380 & 5.386 \\
Open \(0.003\le r<0.005\) & 5.473 & 5.474 \\
Open \(0.001\le r<0.003\) & 5.626 & 5.635 \\
Open \(r<0.001\) & 5.835 & 5.840 \\
\bottomrule
\end{tabular}
\end{center}

Thus \(c=5\) is slightly conservative in the ordinary flat/closed/moderate-open
regime, but not conservative enough for very small open \(r\).  A single global
safe choice for \(10^{-5}\) is \(c\simeq5.85\), but this is unnecessarily early
for most cases.

\subsection{Results for \texorpdfstring{\(10^{-4}\)}{1e-4}}

For \(\varepsilon=10^{-4}\), the corresponding constants are lower:

\begin{center}
\begin{tabular}{lcc}
\toprule
Region & 95th percentile \(c_{\rm eff}\) & max observed \(c_{\rm eff}\) \\
\midrule
Closed \(K=+1\), all tested ratios & 4.189 & 4.193 \\
Open \(K=-1\), \(r\ge0.2\) & 4.195 & 4.204 \\
Open \(0.05\le r<0.2\) & 4.285 & 4.289 \\
Open \(0.02\le r<0.05\) & 4.449 & 4.455 \\
Open \(0.01\le r<0.02\) & 4.590 & 4.594 \\
Open \(0.005\le r<0.01\) & 4.718 & 4.728 \\
Open \(0.003\le r<0.005\) & 4.810 & 4.813 \\
Open \(0.001\le r<0.003\) & 4.943 & 4.962 \\
Open \(r<0.001\) & 5.127 & 5.141 \\
\bottomrule
\end{tabular}
\end{center}

So \(c\simeq4.2\) is the normal-use \(10^{-4}\) value, while a single global
low-open-safe choice is \(c\simeq5.15\).

\begin{figure}[p]
\centering
\includegraphics[width=0.92\linewidth]{c_eff_vs_beta_ratio_eps1em05.png}
\caption{Effective \(c\) values for the \(10^{-5}\)-of-peak threshold as a
function of \(r=\beta/\lambda\).  The closed and moderate-open cases cluster
near \(c\simeq5\), while the low-\(r\) open tail requires larger \(c\).}
\label{fig:ceff1e5}
\end{figure}

\begin{figure}[p]
\centering
\includegraphics[width=0.92\linewidth]{c_eff_vs_beta_ratio_eps0.0001.png}
\caption{Effective \(c\) values for the \(10^{-4}\)-of-peak threshold.  The
normal-use constant is \(c\simeq4.2\), again with a low-\(r\) open correction.}
\label{fig:ceff1e4}
\end{figure}

\section{Practical calibrated constants}

For best practical behaviour use the action equation, Eq.~\eqref{eq:action_rule},
with a constant depending on the target threshold and, only for open models, on
\(r=\beta/\lambda\).

\subsection{Recommended \texorpdfstring{\(10^{-5}\)}{1e-5} rule}

\begin{equation}
\boxed{
 c_{10^{-5}}(K,r)=
 \begin{cases}
 4.95, & K=1,\\
 4.95, & K=-1,\ r\ge0.2,\\
 5.00, & K=-1,\ 0.05\le r<0.2,\\
 4.95+0.42\log_{10}(0.05/r), & K=-1,\ r<0.05,
 \end{cases}
}
\label{eq:c1e5}
\end{equation}
with the cap
\[
  c_{10^{-5}}\le5.85.
\]
This is designed to start near the \(10^{-5}\)-of-peak level without using the
overly conservative global low-\(r\) value in ordinary cases.

\subsection{Recommended \texorpdfstring{\(10^{-4}\)}{1e-4} rule}

\begin{equation}
\boxed{
 c_{10^{-4}}(K,r)=
 \begin{cases}
 4.20, & K=1,\\
 4.21, & K=-1,\ r\ge0.2,\\
 4.30, & K=-1,\ 0.05\le r<0.2,\\
 4.21+0.60\log_{10}(0.05/r), & K=-1,\ r<0.05,
 \end{cases}
}
\label{eq:c1e4}
\end{equation}
with the cap
\[
  c_{10^{-4}}\le5.15.
\]

\section{Algorithm}

Given \(l,K,\beta\), and a threshold \(\varepsilon\):
\begin{enumerate}[label=\arabic*.]
\item Set \(\lambda=l+1/2\) and \(r=\beta/\lambda\).
\item Choose \(c\) from Eq.~\eqref{eq:c1e5} or Eq.~\eqref{eq:c1e4}.
\item Compute the turning point \(\chip\) from \(\SK(\chip)=\lambda/\beta\).
      For \(K=1\), work in the folded interval.
\item Preferably solve the action equation
      \[
      \int_{\chistart}^{\chip}
      \sqrt{\lambda^2/\SK^2(\chi)-\beta^2}\,\dd\chi
      =T_{\lambda,c}.
      \]
\item If speed matters more than exact calibration, use the explicit formula
      \[
      \chistart
      =\max\left[0,
      \chip-
      \frac{c\lambda^{1/3}}
      {\beta(1-K\lambda^2/\beta^2)^{1/6}}
      \right].
      \]
\end{enumerate}

For small \(l\), or whenever \(1-c\lambda^{-2/3}\le0\), the Airy start formula
is no longer meaningful.  In that case use \(\chistart=0\), or determine the
threshold directly from recurrence if avoiding unnecessary integration is
important.

\section{Summary}

The action-inversion recipe is the cleanest generalisation of the flat
\(l-c l^{1/3}\) onset rule:
\[
\boxed{
\int_{\chistart}^{\chip}
\sqrt{\lambda^2/S_K^2(\chi)-\beta^2}\,\dd\chi
= T_{\lambda,c},
\qquad
\lambda=l+\frac12.
}
\]
For \(10^{-5}\)-of-peak accuracy, use \(c\simeq4.95\) in closed and normal open
cases, with the low-open correction in Eq.~\eqref{eq:c1e5}.  For
\(10^{-4}\)-of-peak accuracy, use \(c\simeq4.2\) with Eq.~\eqref{eq:c1e4}.
A single conservative global choice is \(c\simeq5.85\) for \(10^{-5}\) and
\(c\simeq5.15\) for \(10^{-4}\), but these values are unnecessarily early for
most integrations.

\end{document}
